% Paper 8 Figure 2: The AI Engineer's Field Guide to Cellular Sheaves and Cohomology.
% Swiss Modern Systems Graphics specification: visual proofs, hard 90-degree corners,
% metric coordinate drafting grid, and strict chromatic state taxonomy.

\begin{figure}[htbp]
\centering

% =========================================================================
% TOP ROW: Panel (a) and Panel (b)
% =========================================================================
\begin{minipage}[t]{0.485\textwidth}
\vspace{0pt}
\centering
{\normalsize\bfseries (a) What is a Stalk $\mathcal{F}(v)$?}\\[2pt]
{\footnotesize (Agent Private State)}\\[6pt]
\begin{tikzpicture}[
  agentnode/.style={circle, draw=ptDark, fill=ptLightGray, line width=1.1pt, minimum size=1.05cm, inner sep=0pt, font=\sffamily\bfseries\normalsize, text=ptDark},
  fieldbox/.style={rectangle, draw=ptDark, fill=white, line width=1pt, inner sep=6pt, font=\footnotesize, align=left},
  actionbox/.style={rectangle, draw=ptBlue, fill=ptLightGray, line width=1pt, inner sep=5pt, font=\footnotesize, align=center},
]
  % 1. Coordinate Grid Context
  \draw[swissgrid] (-0.25, -0.25) grid (6.85, 3.85);
  \draw[majorgrid] (-0.25, -0.25) grid (6.85, 3.85);

  % Node and stalk
  \node[agentnode] (ag) at (0.6, 2.6) {$v$};
  \node[font=\sffamily\bfseries\scriptsize, text=ptDark, mask] at (0.6, 1.8) {AgentNode};

  \node[fieldbox, text width=4.6cm] (stalk) at (4.2, 2.6) {
    \textbf{\small Vector Stalk} $\mathcal{F}(v) \cong \mathbb{R}^D$ ($x_v \in \mathbb{R}^D$):\\[3pt]
    $\bullet$ \textbf{Epoch Clock}: $t_v \in \mathbb{N}$\\
    $\bullet$ \textbf{Token Spend}: $b_v \in \mathbb{R}_+$\\
    $\bullet$ \textbf{AST Lease}: $h_L \in \mathbb{R}$\\
    $\bullet$ \textbf{Commit Root}: $d_v \in [-1, 1]$
  };
  \draw[ptDark, dashed, thick, ->, shorten >= 3pt] (ag) -- (stalk);

  % Restriction arrow
  \draw[bArr] (4.2, 1.35) -- node[right=3pt, font=\sffamily\bfseries\scriptsize, text=ptBlue, mask] {Restriction $P_{v \trianglelefteq e}$} (4.2, 0.55);

  \node[actionbox, text width=6.4cm, anchor=north] at (3.3, 0.4) {
    \textbf{\small Shared Channel Coordinate Selection}:\\[2pt]
    $P_{v \trianglelefteq e} x_v = [t_v,\, d_v]^T \in \mathcal{F}(e) \cong \mathbb{R}^2$
  };
\end{tikzpicture}
\end{minipage}%
\hfill
\begin{minipage}[t]{0.485\textwidth}
\vspace{0pt}
\centering
{\normalsize\bfseries (b) Cochain Complex $C^0 \xrightarrow{\delta_0} C^1 \xrightarrow{\delta_1} C^2$}\\[6pt]
\begin{tikzpicture}[
  fieldbox/.style={rectangle, draw=ptDark, fill=ptLightGray!50, line width=1pt, inner sep=6pt, font=\footnotesize, align=left, anchor=north},
  actionbox/.style={rectangle, draw=ptBlue, fill=ptBlue!6, line width=1pt, inner sep=5pt, font=\footnotesize, align=center, anchor=north},
]
  % 1. Coordinate Grid Context
  \draw[swissgrid] (-0.25, -0.25) grid (6.85, 3.85);
  \draw[majorgrid] (-0.25, -0.25) grid (6.85, 3.85);

  \node[fieldbox, text width=6.4cm] at (3.3, 3.8) {
    \textbf{\small 0-Cochains ($C^0$)}: Fleet States $x \in \mathbb{R}^{n \cdot D}$\\[3pt]
    \textbf{\small 1-Cochains ($C^1$)}: Channel Differences $g \in \mathbb{R}^{m \cdot d}$\\[3pt]
    \textbf{\small 2-Cochains ($C^2$)}: Triadic Review Curvature $\delta_1 g \in \mathbb{R}^{f \cdot d}$
  };

  \node[actionbox, text width=6.4cm] at (3.3, 1.25) {
    \textbf{\color{ptBlue}\small Fundamental Boundary Nilpotency}:\\[2pt]
    $\delta_1 \circ \delta_0 \equiv 0 \quad \Longleftrightarrow \quad \mathrm{im}(\delta_0) \subseteq \ker(\delta_1)$\\[4pt]
    \textbf{\color{ptBlue}\small Zero False-Alarm Guarantee}:\\[2pt]
    Honest potentials $g = \delta_0 x$ produce exactly\\
    zero triadic review curl ($\delta_1 g \equiv 0$).
  };
\end{tikzpicture}
\end{minipage}

\vspace{0.35cm}
\hrule height 0.8pt \relax
\vspace{0.35cm}

% =========================================================================
% BOTTOM ROW: Panel (c) and Panel (d)
% =========================================================================
\begin{minipage}[t]{0.485\textwidth}
\vspace{0pt}
\centering
{\normalsize\bfseries (c) Gossip Trees vs.\ Closed Cycles}\\[2pt]
{\footnotesize (Acyclic Silence vs.\ Certified Traps)}\\[6pt]
\begin{tikzpicture}[
  agentnode/.style={circle, draw=ptDark, fill=ptLightGray, line width=1pt, minimum size=7.2mm, inner sep=0pt, font=\sffamily\bfseries\footnotesize, text=ptDark},
  roguenode/.style={circle, draw=ptRed, fill=ptRed!20, line width=1.2pt, minimum size=7.2mm, inner sep=0pt, font=\sffamily\bfseries\footnotesize, text=ptRed},
  actionbox/.style={rectangle, line width=1pt, inner sep=5pt, font=\footnotesize, align=center, anchor=north},
]
  % 1. Coordinate Grid Context
  \draw[swissgrid] (-0.25, -0.75) grid (6.85, 4.45);
  \draw[majorgrid] (-0.25, -0.75) grid (6.85, 4.45);

  % Left: Open Tree
  \node[font=\sffamily\bfseries\scriptsize, text=ptDark, align=center, mask] at (1.6, 4.2) {Open Delegation Tree\\[1pt]{\scriptsize (LangChain / CrewAI)}};
  \node[agentnode] (t0) at (1.6, 3.2) {$v_0$};
  \node[roguenode] (t1) at (0.8, 2.2) {$v_1^*$};
  \node[agentnode] (t2) at (2.4, 2.2) {$v_2$};
  \node[agentnode] (t3) at (0.4, 1.3) {$v_3$};
  \node[agentnode] (t4) at (1.2, 1.3) {$v_4$};
  \node[agentnode] (t5) at (2.4, 1.3) {$v_5$};

  \draw[dArr] (t0) -- (t1);
  \draw[dArr] (t0) -- (t2);
  \draw[dArr] (t1) -- (t3);
  \draw[dArr] (t1) -- (t4);
  \draw[dArr] (t2) -- (t5);

  \node[actionbox, draw=ptRed, fill=ptRed!5, text width=3.1cm] at (1.6, 0.4) {
    \textbf{\color{ptRed} $\ker(\delta_0^T) = \{0\}$ (Silent!)}\\[2pt]
    Acyclic trees swallow lies into scalar potentials; $r \equiv 0$.
  };

  % Right: Closed Cycle
  \node[font=\sffamily\bfseries\scriptsize, text=ptDark, align=center, mask] at (5.3, 4.2) {Closed Review Cycle\\[1pt]{\scriptsize (Autonomous Control Plane)}};
  \node[agentnode] (c0) at (4.3, 3.1) {$u_0$};
  \node[roguenode] (c1) at (6.3, 3.1) {$u_1^*$};
  \node[agentnode] (c2) at (6.3, 1.3) {$u_2$};
  \node[agentnode] (c3) at (4.3, 1.3) {$u_3$};

  \draw[dArr] (c0) -- (c1);
  \draw[rArr] (c1) -- node[right=2pt, font=\sffamily\bfseries\scriptsize, text=ptRed, mask] {Lie} (c2);
  \draw[dArr] (c2) -- (c3);
  \draw[dArr] (c3) -- (c0);

  \draw[rArr, dashed] (5.3, 2.2) arc (30:330:0.42);
  \node[font=\sffamily\bfseries\small, text=ptRed, mask] at (5.3, 2.2) {$\rho \neq 0$};

  \node[actionbox, draw=ptBlue, fill=ptBlue!6, text width=3.1cm] at (5.3, 0.4) {
    \textbf{\color{ptBlue} $\ker(\delta_0^T) \neq \{0\}$ (Trapped!)}\\[2pt]
    Cycle traps contradiction with certified residual $r > 0$.
  };
\end{tikzpicture}
\end{minipage}%
\hfill
\begin{minipage}[t]{0.485\textwidth}
\vspace{0pt}
\centering
{\normalsize\bfseries (d) Gauge Skew vs.\ Triadic Curl}\\[6pt]
\begin{tikzpicture}[
  gaugenode/.style={circle, draw=ptBlue, fill=ptBlue!12, line width=1.1pt, minimum size=7.5mm, inner sep=0pt, font=\sffamily\bfseries\footnotesize, text=ptBlue},
  roguenode/.style={circle, draw=ptRed, fill=ptRed!20, line width=1.2pt, minimum size=7.5mm, inner sep=0pt, font=\sffamily\bfseries\footnotesize, text=ptRed},
  actionbox/.style={rectangle, line width=1pt, inner sep=5pt, font=\footnotesize, align=center, anchor=north},
]
  % 1. Coordinate Grid Context
  \draw[swissgrid] (-0.25, -0.75) grid (6.85, 4.45);
  \draw[majorgrid] (-0.25, -0.75) grid (6.85, 4.45);

  % Left: Gauge Skew
  \node[font=\sffamily\bfseries\scriptsize, text=ptBlue, align=center, mask] at (1.6, 4.2) {Gauge Skew\\[1pt]{\scriptsize (Benign Turn Lag)}};
  \node[gaugenode] (g0) at (0.6, 1.3) {$t_0$};
  \node[gaugenode] (g1) at (2.6, 1.3) {$t_1$};
  \node[gaugenode] (g2) at (1.6, 3.4) {$t_2$};

  \draw[bArr] (g0) -- (g1);
  \draw[bArr] (g1) -- (g2);
  \draw[bArr] (g0) -- (g2);
  \node[font=\sffamily\bfseries\small, text=ptBlue, mask] at (1.6, 2.1) {$\delta_1 g = 0$};
  \node[font=\sffamily\bfseries\scriptsize, text=ptBlue, mask] at (1.6, 1.7) {Curl-free};

  \node[actionbox, draw=ptDark, fill=ptLightGray, text width=3.1cm] at (1.6, 0.4) {
    \textbf{\color{ptDark} Action: Do Nothing!}\\[2pt]
    Benign step latency;\\
    self-resolves next turn.
  };

  % Right: Triadic Curl
  \node[font=\sffamily\bfseries\scriptsize, text=ptRed, align=center, mask] at (5.3, 4.2) {Triadic Curl\\[1pt]{\scriptsize (Review Bug / Paradox)}};
  \node[roguenode] (r0) at (4.2, 1.3) {Coder};
  \node[roguenode] (r1) at (6.4, 1.3) {Critic};
  \node[roguenode] (r2) at (5.3, 3.4) {Mgr};

  \draw[rArr] (r0) -- (r1);
  \draw[rArr] (r1) -- (r2);
  \draw[rArr] (r2) -- (r0);

  \draw[rArr, dashed] (5.3, 2.1) arc (30:330:0.38);
  \node[font=\sffamily\bfseries\small, text=ptRed, mask] at (5.3, 2.1) {$\delta_1 g \neq 0$};

  \node[actionbox, draw=ptRed, fill=ptRed!5, text width=3.1cm] at (5.3, 0.4) {
    \textbf{\color{ptRed} Action: Re-prompt!}\\[2pt]
    Review contract violated;\\
    restart critique loop.
  };
\end{tikzpicture}
\end{minipage}

\vspace{0.35cm}
\begin{tikzpicture}
  \node[rectangle, draw=ptDark, fill=ptLightGray, line width=1pt, text width=0.96\linewidth, inner sep=7pt, align=center] {
    {\normalsize\textbf{Swarm Legibility Ratio}: \quad
    $\displaystyle \mathcal{L}(g) = \frac{\|h\|_2^2}{\|h\|_2^2 + \|\delta_1^* \psi\|_2^2} \in [0, 1]$}\\[4pt]
    {\small $\mathcal{L}(g) \approx 0 \implies$ \textbf{\color{ptRed}Local Review Hallucination} (Re-prompt critique triad) \quad $\vert$ \quad $\mathcal{L}(g) \approx 1 \implies$ \textbf{\color{ptBlue}Macro Network Split} (Harmonic partition cavity)}
  };
\end{tikzpicture}

\caption{The AI Engineer's Field Guide to Cellular Sheaves and Cohomology. (a) An agent's private state vector is modeled as a vector stalk $\mathcal{F}(v)$ with coordinate restrictions. (b) The cochain complex $C^0 \to C^1 \to C^2$ satisfies $\delta_1 \circ \delta_0 = 0$, guaranteeing zero false alarms for honest swarms. (c) Hierarchical delegation trees have $\ker(\delta_0^T) = \{0\}$ and silently swallow lies, whereas closed review cycles trap contradictions with non-zero completion residuals $r > 0$. (d) Discrete Hodge theory separates benign turn latency (gauge processes) from review hallucinations (triadic curl) and macro-network cavities.}
\label{fig:fieldguide}
\end{figure}
